\section{Public-Model Audit}
\label{sec:audit}

We apply SNF-Bench to publicly released long-horizon generators. The goal is not a universal leaderboard but a determination of how checkpoints occupy the static-fidelity/flow-persistence space under a single recorded T2V setting.

\paragraph{Systems and setting.}
The audit covers \SNFnPublicTTV{} publicly released autoregressive or streaming generators with accessible checkpoints and reproducible inference code, frozen as a set before the sweep. Every T2V checkpoint runs through one common inference configuration rather than each method's released defaults: a four-step denoising schedule, guidance $5.0$, six frames per block and a fixed seed, with only the checkpoint and horizon varying. All checkpoints are evaluated at the same output resolution, frame rate, target horizons, and wrapper settings. Accordingly, the results isolate checkpoint behavior under this common deployment rather than reconstructing each method's published inference procedure. A checkpoint distilled for a different step budget is therefore not observed at its own operating point. The I2V track separately distinguishes recorded native outputs from checkpoints run through its common long-horizon wrapper (Sec.~\ref{sec:results_i2v}).

\paragraph{Controls and statistics.}
Within each item every system receives identical conditioning, and the seed set is fixed before the sweep. Seed control guarantees reproducibility within a system, but because implementations consume randomness differently, identical seeds do not induce matched noise draws across systems; cross-system comparison therefore rests on the paired-item design. Videos are retained in their recorded form and the normalization of Sec.~\ref{sec:metrics} applied unchanged. Scores are computed per item and averaged over items; the sensitivity of every conclusion to category macro-averaging is reported in the supplementary material. Uncertainty is estimated by percentile bootstrap~\cite{efron1979bootstrap} over items, and we report $95\%$ marginal intervals. These intervals summarize per-checkpoint uncertainty; all headline separation claims use paired bootstrap distributions of per-prompt differences, reported in Supplementary Table~S3. Operational limits---unsupported horizons, failed generations, memory exhaustion, invalid outputs---are reported as rates against the frozen manifest and never removed after inspection. To test whether the decomposition changes interpretation, the same outputs are additionally scored with established whole-frame metrics and the induced orderings compared. The focus is interpretation change, not a claim that existing metrics are wrong.
